\section{A uniform certificate through a moving collision}
\label{sec:v18-microscopic}
We now use all three parts of the argument on one family. The finite marked
reference law and the uniform scattering geometry are those of
Theorem~\ref{thm:v17-markedgap} and Lemma~\ref{lem:v17-scattering}.
They are proved above, including the no-collision control. We use a
continuous observation of the same detector so that finite reports and
conditional processes are generated by one experiment.

Let $z$ denote a preparation in the boxes of
\eqref{eq:v17-packets}, and let $\nu$ be that prepared law. The sign
of the impact orientation is $2B-1$, and the conserved center-of-mass
longitudinal sign is the actual mark $W$. Their correlation is prepared:
$\Pp(B=b,W=w)=5/16$ for $b=w$ and $3/16$ otherwise. The contact diameter
$d$ ranges over $I=[9/10,11/10]$. Let $X_d(t,z)$ be the two-particle phase
path, with outgoing velocities used at the collision, and define
\[
 H_d(t,z)=h(v_{1,y}^d(t,z)),\qquad
 Y_t^d=\int_0^tH_d(s,z)ds+\sigma B_t^d,
 \quad 0<\sigma\le1/24.
\]
The Brownian noise is independent of the preparation. The source experiment
has the same physical preparation and mark but constant reports. The target
reports at times $2$ and $3$ are
\begin{equation}\label{eq:v18-windowreports}
 Z_1^d=\1\{2(Y_2^d-Y_{3/2}^d)>0\},\qquad
 Z_2^d=A\1\{2(Y_3^d-Y_{5/2}^d)>0\}.
\end{equation}
The acquisition gate $A\in\{0,1\}$ is chosen after $Z_1^d$; it does not
act on the mechanical trajectory. The continuous path is the auxiliary
observation experiment for conditional consumers; the finite simulator is
not given this path or an uncharged buffer for computing its increments.
The target sensor implements the declared window acquisition.

\begin{lemma}[Quantitative change of the microscopic path]
\label{lem:v18-movingcollision}
For all preparations in the specified boxes, the collision time $t_c(d)$
and outgoing tagged velocities obey
\[
 |t_c(d')-t_c(d)|\le|d'-d|,
 \qquad |v_i^+(d')-v_i^+(d)|\le30|d'-d|\quad(i=1,2).
\]
Writing $\Delta=|d'-d|$, one has
\begin{equation}\label{eq:v18-signalmodulus}
 \sup_z\int_0^3|H_{d'}(t,z)-H_d(t,z)|^2dt
                         \le\Delta+2700\Delta^2.
\end{equation}
The phase paths converge uniformly over the preparation in the
Skorokhod $J_1$ metric as $d'\to d$. More precisely, for the product
metric taking the maximum of the per-particle Euclidean position and
velocity distances, a time change aligning the collision times gives
$J_1$ distance at most $100\Delta$.
\end{lemma}
\begin{proof}
Write $r$ and $g$ for the incoming relative position and velocity, and
$n(d)=(r+t_c(d)g)/d$ for the contact normal. The collision root satisfies
$|r+t_cg|=d$. The nongrazing estimate from
Lemma~\ref{lem:v17-scattering} is
$-g\cdot n>1443/1000$. Differentiating the root with respect to $d$ gives
\[
 t_c'(d)=\frac1{g\cdot n(d)},\qquad |t_c'(d)|<1.
\]
The incoming bound is $|g|<3$, so
\[
 |n'(d)|\le\frac{|t_c'(d)|\,|g|+1}{d}<5.
\]
Since $v_1^+=V+g/2-(g\cdot n)n$ and
$v_2^+=V-g/2+(g\cdot n)n$, differentiation bounds either velocity
derivative by $2|g|\,|n'|<30$.

Before collision the detector is zero, and after collision it is a
constant of absolute value at most one. Its Lipschitz constant is one.
The two signals therefore differ on the interval between the two collision
times by at most one and afterwards by at most $30\Delta$. The intervening
interval has length at most $\Delta$, and the whole observation lasts
three units. This proves \eqref{eq:v18-signalmodulus}.

Choose the increasing piecewise affine time change fixing $0,3$ and sending
$t_c(d)$ to $t_c(d')$. Its uniform displacement is at most $\Delta$.
Aligned velocities differ by at most $30\Delta$, with the same outgoing
convention at the jump. Each particle's speed is below four and its
velocity jump has norm below three. Without time change, the position
difference is bounded by $3\Delta+3(30\Delta)=93\Delta$, by integrating
the velocity difference before and after the two jump times. Aligning time
adds at most $4\Delta$. The asserted bound follows, with a harmless
rounding to $100\Delta$.
\end{proof}

\begin{theorem}[Microscopic-to-audit and conditional-law composition]
\label{thm:v18-microscopic}
For the single family \eqref{eq:v18-windowreports}, all of the following
hold.

\smallskip\noindent
\textup{(a)} At the same width-one private register,
\[
 \dpr>9/20,\qquad \dvi_r>9/20\quad(r\ge1),\qquad
 \dhi_2<13/100,
\]
uniformly for $d\in I$ and $0<\sigma\le1/24$.
When $d\le2/5$ the detector is zero on $[0,3]$, and all three
finite-experiment deficiencies are zero.

\smallskip\noindent
\textup{(b)} There is one explicit rational reset-audit rule with one
all-on training row, eight marked outcomes and $n=2800$ training runs,
whose expected validation score exceeds $2/5$ against every private
width-one simulator throughout that same parameter region. In detail,
for counts $C$ choose an event $A_C\subseteq\{0,1\}^3$ maximizing
\begin{equation}\label{eq:v18-concreteaudit}
 Q_*(A_C)-C(A_C)/2800,
\end{equation}
with lexicographic tie breaking. Here $Q_*$ is the rational marked
reference law of $(B,B,W)$. The independent validation uses the actual
physical target, not $Q_*$. This finite rule has at most
$\binom{2807}{7}$ training-count states at any scheduled epoch; a
conservative complete auditor bound is
$8(256+1)\binom{2807}{7}$. The state cost belongs to the auditor.
The certificate is an expected-score inequality, not a confidence claim
about one realized score.

\smallskip\noindent
\textup{(c)} Fix $\sigma>0$ and let $d_n,d\in I$ with $d_n\to d$.
Set
\begin{equation}\label{eq:v18-physicalepsilon}
 \epsilon_n=\min\left\{1,
       \frac{\sqrt{|d_n-d|+2700|d_n-d|^2}}{2\sigma}\right\}.
\end{equation}
The joint preparation--continuous-observation laws differ by at most
$\epsilon_n$. Every finite resource deficiency above changes by at most
$\epsilon_n$, and a fixed audit changes by at most $\epsilon_n$ because
its source behavior is unchanged. On a latent-preserving coupling,
$X_{d_n}\to X_d$ in $J_1$ and each bounded posterior converges in
supremum $L^2$ with bound $84\epsilon_n$. The full and projected
likelihoods converge in reference supremum $L^2$, and jointly in probability
with the microscopic paths on that coupling. The projected stochastic
exponentials, bounded entropic backward values, localized integrands and
stopping consumers have the conclusions of
Proposition~\ref{prop:v18-innovation},
Theorem~\ref{thm:v18-entropic}, and Corollary~\ref{cor:v18-stopping}.

\smallskip\noindent
\textup{(d)} On compact parameter sets with $\sigma\ge\sigma_0>0$, the
finite marked probabilities have uniform rational enclosures. The
observable chart, strict same-budget polynomial verification and rational
upper tables are therefore effective in the exact sense of
Propositions~\ref{prop:v18-chart} and \ref{prop:v17-effective}.
\end{theorem}
\begin{proof}
Every collision occurs before time $3/2$. Conditional on preparation, the
two normalized window increments in \eqref{eq:v18-windowreports} are
independent Gaussians with mean $h(v_{1,y}^+)$ and standard deviation
$\sqrt2\sigma\le1/12$. The outgoing detector has sign $2B-1$ and absolute
value greater than $1/4$. Coupling each noisy report to $B$ gives total
marked variation less than $1/300$ from the exact finite marked experiment,
uniformly for every acquisition gate. The finite calculation of
Theorem~\ref{thm:v17-markedgap} gives
\[
 d_* =\sqrt{5/2}-9/8,\qquad
 \delta_*^{\rm p}=\delta_*^{\rm v}=d_*,\quad
 \delta_{*,2}^{\rm h}=1/8.
\]
Thus private and visible values exceed $d_*-1/300>9/20$, and the hidden
value is less than $1/8+1/300<13/100$. The two-valued hidden witness stores
its selector, not an observed mark. In the no-collision regime both window
reports before gating are independent fair noise, independent of $W$.
Fresh private coins reproduce them exactly, including feedback gates.
This proves (a).

For (b), use only the all-on feedback row. Any private width-one simulator
in this row emits independent bits with probabilities $p,q$, independently
of the actual $W$. The calculation in
Theorem~\ref{thm:v17-markedgap} already gives minimum variation $d_*$ for
this row. Its alphabet $(z_1,z_2,W)$ has eight elements, hence $D=7$.
The explicit empirical audit in Theorem~\ref{thm:v18-audit}, trained using
$Q_*$, has expected ideal score at least
\[
 d_*-\sqrt{7/2800}=d_*-1/20.
\]
The candidate training and validation laws do not change when the target
is replaced by the physical target. Conditional on any selected event,
its target probability changes by less than $1/300$. Therefore the actual
score is strictly greater than
\begin{equation}\label{eq:v18-uniformcertificate}
 d_*-\frac1{20}-\frac1{300}>\frac25.
\end{equation}
For completeness, the last strict inequality is the positive-root
comparison
$\sqrt{5/2}>947/600$; its squared margin is
$5/2-(947/600)^2=3191/360000>0$.
All probabilities in $Q_*$ are $0$, $3/16$ or $5/16$, so
\eqref{eq:v18-concreteaudit} is a finite rational rule independent of
$d,\sigma,p,q$. At training epoch $k$ the number of eight-category count
vectors is $\binom{k+7}{7}\le\binom{2807}{7}$. A buffer of eight states
records the current two reports and terminal mark; 256 event masks and
an idle symbol suffice during event selection and validation. The
conservative product bound follows. The clock and the source-reset
interface are the declared resources. In particular the proof does not
enumerate or claim to store a small complete coefficient table.

For (c), Lemma~\ref{lem:v18-movingcollision} gives the signal-distance
hypothesis of Theorem~\ref{thm:v18-gaussian}, with $T=3$ and $M=1$.
It also supplies its microscopic $J_1$ hypothesis. That theorem gives
\eqref{eq:v18-physicalepsilon} and all the joint process assertions.
The finite acquisition is a common measurable causal map of the continuous
observation, so marked variation contracts for every gate. The source is
blind and independent of $d$. Consequently its executable candidate
behaviors are identical across $d$, and
Proposition~\ref{prop:v18-audittransport} has $\alpha=0$ and
$\beta=\epsilon_n$. The same same-budget comparison applies in each
visible fiber and in each hidden mixture. The exponential and backward
claims now follow from the identified projected density and the explicitly
bounded entropic consumer. This is where the finite certificate and the
conditional-law theorem meet the same microscopic approximation.

Finally, the collision root is uniformly separated from grazing, its
formula is algebraic on the rational preparation boxes, and the normalized
window noise has variance $2\sigma^2$ bounded away from zero on the stated
compact sets. The box subdivision and Gaussian-tail enclosure argument
in Proposition~\ref{prop:v17-collisioncompute} therefore applies with
that variance. It returns rational enclosures with a computable uniform
error. The rational quotient algorithm and same-budget strict certificate
procedure then prove (d). Uniformity as $\sigma\downarrow0$ is not asserted.
\end{proof}

The all-on audit in part (b) is a single witness for the lower bound in the
full feedback problem: a maximum over feedback trees includes this row.
The visible lower bound is instead the exact fiber theorem; no claim is
made that an auditor with a seed-independent state bound can store an
arbitrarily large visible-seed transcript. This separates the two resource
budgets rather than silently moving cost from one to the other.

The finite certificate is also preserved under an additional rational
presentation error $\rho$: its lower score decreases by at most $\rho$
when only the target is changed, and by the reset-cost estimate when the
source is changed. Thus the certificate is an effective consumer of the
physical presentation theorem, not an isolated exact finite example.

Theorem~\ref{thm:v18-microscopic} and the preceding results prove
Theorem~\ref{thm:v18-spine}. The C2-type consequence is a joint
changing-trajectory/likelihood/posterior theorem with an identified
stochastic exponential and a fully specified bounded entropic backward
consumer. Nothing in that conclusion identifies geometric similarity with
kinetic scaling, or supplies particle-uniform recollision estimates by
renaming this fixed-particle model.
